\solution{Жұлдыз тұрақтылығы}{10.0}

\tasksol Геометриялық тұрғыдан алғанда, Жер $2R_E$ қашықтығына ығысқан жарты жыл ішінде $\Delta\varphi$ бұрышы еңсеріледі. Тригонометрия бойынша

\begin{eq}[points=0.2]
d = R_E \cot\left(\frac{\Delta\varphi}2\right)\approx\frac{2R_E}{\Delta\varphi}.
\end{eq}

Сандық түрде,

\begin{eq}[points=0.2]
d=\frac{2\times\num{1.5e11}}{\num{66.7e-3}\times\frac1{3600}\times\frac\pi{180}}=\qty{9.28e17}{\m}=\qty{98.0}{\ly}
\end{eq}

\criterion[type=penalty, points=-0.1]{Жарық жылы (\unit{\ly}) бірліктерімен жазылмаған}

\tasksol Бізде барлығы

\begin{eq}[points=0.2]
N = n\cdot \frac43\pi R^3
\end{eq}

молекула бар. Әр молекуланың энергиясы $\frac32 kT$ тең (нөлге жуық температуралар үшін біз тек ілгерілемелі еркіндік дәрежелерін ескереміз). Сондықтан

\begin{eq}[points=0.2]
U = \frac32 NkT.
\end{eq}

Осы екі теңдеуді біріктіріп, келесі нәтижені аламыз.

\tasksol Бұл есепте біртекті шар ішіндегі гравитациялық өріс кернеулігі радиусқа байланысты сызықты түрде өсетінін еске түсіру жеткілікті (мұны гравитациялық өріс үшін Гаусс теоремасы немесе Ньютонның сфералық қабықша туралы теоремасы арқылы көрсетуге болады). Бетінде

\begin{eq}[points=0.2]
g_0=\frac{GM}{R^2},
\end{eq}

болса, біз келесі функцияны аламыз:

\begin{eq}[points=0.3]
g(r)= g_0\cdot\frac rR.
\end{eq}

\tasksol Бұл есепке екі түрлі жолмен қарауға болады. Жылдамырақ жолы --- $|E_G|$ шамасы бұлтты шексіздікке дейін ``апару'' үшін, яғни тартылыс энергиясы нөлге тең болуы үшін тартылыс күштеріне қарсы атқарылуы тиіс минималды жұмыстың мағынасын беретінін түсіну\footnote{Жалпы жағдайда, егер біз электростатикалық есеппен жұмыс істесек, $\max\{0, -E_G\}$ деп айтар едік.}. Айталық, бұлт айнымалы $r$ радиусына дейін кішірейтілді. Онда кіші сфералық $dm$ элементінің потенциалдық энергиясы

\begin{eq}[points=0.3]
\phi(r) = -\frac{GM(r)}r=-\frac{GM r^2}{R^3}.
\end{eq}

Онда $E_G$ келесі интегралға тең:

\begin{eq}[points=0.2]
E_G = \int_0^R \phi(r)\,dm.
\end{eq}

Баламалы, бірақ сәл егжей-тегжейлі әдіс ретінде шар ішіндегі гравитациялық потенциалдың $\varphi(r)$ (яғни масса бірлігіне шаққандағы потенциалдық энергия) таралуын қарастырып, содан кейін бұлттың барлық элементтері бойынша қосындылауға болады. Бұл әдіспен алатынымыз\footnote{Мұнда мұқият болу керек. Дәлірек айтсақ, біз $\varphi = -\frac{\partial g_r}{\partial r}$ деп анықтаймыз, бірақ сфералық координаттар жүйесінде радиалды компонента бұлттан ``сыртқа'', ал $g(r)$ --- ``ішке'' бағытталғанын ескерсек, $g_r=-g(r)$ болады.}
$$
\varphi(r) - \varphi(R) = \int_R^r g(r)\, dr.
$$
Мұнда
$$
\varphi(R) = -\frac{GM}R,
$$
сондықтан \solref{3}-тегі нәтижені қолдана отырып,
$$
\varphi(r) = -\frac{GM}R\ab(\frac32-\frac{r^2}{R^2}).
$$
Потенциалдық энергия интегралдау арқылы табылады
$$
E_G = \frac12\int_0^R \varphi(r)\,dm,
$$
мұнда $dm=\rho\cdot 4\pi r^2\,dr= 3M r^2\,dr/R^3$, ал $1/2$ коэффициенті екі рет есептелген өзара әрекеттесулерді ескереді. Интегралдай отырып, келесі нәтижені аламыз:

\begin{eq}[points=0.5]
E_G=-\frac35\frac{GM}R\iff f=\frac35.
\end{eq}

\tasksol Келесі қатынасты қосу керек

\begin{eq}[points=0.2]
M = mN,
\end{eq}

мұндағы $N$ \eqref{3}-те өрнектелген. Толық энергия $E=E_G+U$ тең, сондықтан критикалық жағдайда $|E_G|=U$. \eqref{3}--\eqref{4} және \eqref{9} нәтижелерін біріктіре отырып, келесіні аламыз:

\begin{eq}[points=0.3]
    M_J=\frac{5k T}{4Gm^2}\sqrt\frac{15 k T}{2\pi nG}=\qty{2.60e31}{\kg}.
\end{eq}

\criterion[type=remark]{$M_J$-ны кг-мен жазу міндетті емес}

Күн массасы бірліктерімен,
\begin{eq}[points=0.1]
M_J=13 M_\odot.
\end{eq}

Әрине, бұл --- \textit{минималды} масса. Басқа параметрлер бекітілген жағдайда, кез келген $M>M_J$ үшін коллапс орын алады.

\tasksol Массасы $dM=\rho(r)\cdot S dr$ болатын дискті қарастырайық. Оған жоғарыдан $P(r+dr)S$ қысым күші, ал төменнен $P(r)S$ әсер етеді. Сонымен қатар, төмен бағытталған $dM\cdot g(r)$ күші бар, соның салдарынан қысым градиенті пайда болады. Гидростатикалық тепе-теңдік теңдеуі

\begin{eq}[points=0.2]
P(r)S - P(r+dr)S-dM\cdot g(r)=0.
\end{eq}

Енді $g(r)$-ді өрнектеу қалды. \eqref{6}-ны жалпылап, жалпы түрде былай жазамыз:

\begin{eq}[points=0.2]
g(r)=\frac{GM(r)}{r^2}.
\end{eq}

$P(r+dr)-P(r)=dP$ екенін ескерсек, аламыз

\begin{eq}[points=0.3]
\frac{dP}{dr}=-\frac{GM(r)\rho(r)}{r^2}.
\end{eq}

\tasksol Техникалық тұрғыдан, $\rho(r)=\operatorname{const}$ деп алмай, $E_G$ формуласын жалпы түрде қайта қорытып шығару керек, бұл \eqref{9}-дан өзгеше болады. $M(r)$-ді ашпай, \eqref{7}-ні қолданамыз. Онда біз $dm=\rho(r)\cdot4\pi r^2\,dr$ деп жазатын \eqref{8} интегралын аламыз. Онда $r\cdot dP/dr=\phi(r)\rho(r)=\phi(r)\,dm/dV$, немесе

\begin{eq}[points=0.4]
E_G = \int_0^{V_\odot} r\frac{dP}{dr}\,dV=\int_0^{R_\odot}4\pi r^3\,dP=\int_0^{R_\odot}-4\pi G M(r)\rho(r) r\,dr.
\end{eq}

\criterion[type=remark]{Егер $\rho(r)$ және $P(r)$ жалпы түрде сақталса, кез келген ұқсас жазба қабылданады}

Қатысушының қалауына қарай бұл интегралды түрліше жазуға болады: жоғарыда үш мүмкін мысал келтірілген және $\rho(r)=\operatorname{const}$ болжамы қолданылмаған кез келген дұрыс интеграл жазбасы үшін балл беріледі. Әрі қарай бұл былай жазылады:

\begin{eq}[points=0.3]
E_G=\int 3V\,dP=3PV\Big|_0^{R_\odot}-\int_0^{V_\odot}3P\,dV=-3\int_0^{V_\odot}P\,dV.
\end{eq}

\criterion[type=remark]{Көрсетілген бөліктеп интегралдау тек $PV\big|_0^{r_\odot}=0$ екендігі көрсетілген жағдайда \textbf{ғана} қабылданады}

мұнда бірінші `=' таңбасында бөліктеп интегралдау қолданылды, сонымен қатар $PV|_0^{R_\odot}$ мүшесі төменгі шекте көлемнің нөлге теңдігінен, ал жоғарғы шекте қысымның нөлге теңдігінен жойылып кетеді. Нәтижесінде

\begin{eq}[points=0.3]
E_G = -3\langle P\rangle V_\odot\iff \beta=-3.
\end{eq}

\tasksol \eqref{9}-ға $f=1$ мәнін қойып және оны \eqref{18}-ге теңестіре отырып, бірден келесіні аламыз:

\begin{eq}[points=0.3, numpoints=0.2]
\langle P\rangle = \frac{GM_\odot^2}{4\pi r^4_\odot}=\qty{9e13}{\Pa}.
\end{eq}

\tasksol Идеал газ күйінің теңдеуі келесі түрде жазылады:

\begin{eq}[points=0.2]
    \langle P\rangle = \frac{\langle\rho\rangle}{\overline m}k T_0 \iff \langle P\rangle V_\odot = \langle\nu\rangle R T_0,
\end{eq}

мұнда $\langle\rho\rangle=\frac{M_\odot}{V_\odot}\iff\langle\nu\rangle=\frac{M_\odot}{\overline m N_A}$. Нәтижесінде

\begin{eq}[points=0.2, numpoints=0.1]
T_0=\frac{\overline m \langle P\rangle V_\odot}{k M_\odot}= \frac{GM_\odot\overline m}{3kR_\odot}=\qty{4.7e6}{\K}.
\end{eq}

\tasksol Стефан--Больцман заңын стандартты қолдану. Былай жазып,

\begin{eq}[points=0.2]
    L_\odot=\sigma T^4_s\cdot4\pi R_\odot^2,
\end{eq}

табатынымыз

\begin{eq}[points=0.2, prefix={Сандық мәні}]
T_s=\ab(\frac{L_\odot}{4\pi\sigma R_\odot^2})^{1/4}\approx\qty{5800}{\K}.
\end{eq}

\tasksol Фотонның қорытқы орын ауыстыруы

\begin{eq}[points=0.2]
    \vec D = \vec l_1+\vec l_2+\dots+\vec l_N.
\end{eq}

Фотонның қозғалысы таза стохастикалық (кездейсоқ) деп жуықталғандықтан, қозғалыстың барлық бағыттары бірдей ықтималды, сондықтан орташа мәні

\begin{eq}[points=0.3]
    \langle\vec D\rangle=0.
\end{eq}

\tasksol \eqref{24} өрнегінің квадратын алайық. Шығатыны

\begin{eq}[points=0.3, override={D^2 = l_1^2+\dots+l_N^2+2\ab(\vec l_1\cdot\vec l_2+\vec l_1\cdot\vec l_3+\dots)}]
    D^2 = l_1^2+\dots+l_N^2+2\ab(\vec l_1\cdot\vec l_2+\vec l_1\cdot\vec l_3+\dots)=\sum_{i=1}^N l_i^2 + 2\sum_{i\neq j}\vec l_i\cdot\vec l_j,
\end{eq}

\criterion[type=penalty, points=-0.2]{Түсіндірмесіз жұптық қосындылар көрсетілмеген}

мұнда жақша ішіндегі өрнек әртүрлі индекстер бойынша барлық жұптық қосындыларды алады. $|l_i|=l\ \forall i$ екенін ескерсек, бірінші қосынды $Nl^2$ береді. Фотон қозғалысының стохастикалықтығын тағы да ескерсек, орташа алғанда екінші қосынды нөлге тең болатынын көреміз. Нәтижесінде

\begin{eq}[points=0.2]
    \langle D^2\rangle = Nl^2.
\end{eq}

Күннен шыққан фотонның орташа квадраттық ығысуы шамамен $r^2_\odot$ тең болуы керек. Сонымен, фотонға Күннен шығу үшін шамамен

\begin{eq}[points=0.2]
    \langle D^2\rangle\approx R_\odot^2\iff N=\ab(\frac{R_\odot}l)^2
\end{eq}

қадам қажет, ал бұл келесі уақытқа сәйкес келеді:

\begin{eq}[points=0.1]
    t = N\cdot\frac lc.
\end{eq}

Сонымен,

\begin{eq}[points=0.2]
    t=\frac{R_\odot^2}{cl}.
\end{eq}

\tasksol Шартта сипатталғандай, біз былай жаза аламыз:

\begin{eq}[points=0.3]
    L_\odot' = \sigma T_0^4\cdot 4\pi R_\odot^2.
\end{eq}

Бұны \eqref{22}-мен және шарттағы ақпаратпен байланыстыра отырып, аламыз

\begin{eq}[points=0.2, numpoints=0.2]
l = R_\odot\ab(\frac{T_s}{T_0})^4\approx\qty{1.67}{\mm}.
\end{eq}

Бұл \eqref{30}-дан уақыттың сандық мәні келесіге тең екенін көрсетеді:

\begin{eq}[points=0.3, prefix={Сандық мәні}]
    t=\qty{9.67e11}{\s}\approx\num{30600}~\text{жыл}.
\end{eq}

\tasksol \eqref{21} теңдеуін алып, оны \eqref{31} өрнегіне қоямыз. \eqref{31}-ді $l/R_\odot$-ге көбейтсек, бұл $L_\odot$ үшін өрнек болады. Онда шығатыны

\begin{eq}[points=0.3]
    L_\odot = \frac{4\pi \sigma G^4 m^4 l}{81 k^4}\cdot\frac{M_\odot^4}{R_\odot^3}.
\end{eq}

Осы өрнектен $\langle\rho\rangle$-ны шығара отырып, қорытамыз

\begin{eq}[points=0.2]
    C=\frac{16 \pi^2 \sigma G^4 m^4 l}{243 k^4},
\end{eq}

\begin{eq}[points=0.2]
n=3.
\end{eq}

\tasksol Соңғы теңдеуден бізде

\begin{equation}
\log L = \log L + n\cdot \log M,
\end{equation}

сондықтан $\log L$--$\log M$ координаттарында біз шарт бойынша графикте берілген нүктелерді түзу сызықпен жуықтаймыз. Икс осі логарифмдік масштабқа ие болса да, логарифм алып, осьтегі мәндерді қайта есептеу керек екеніне назар аудару қажет. Онда төмендегі суреттегідей нәтиже аламыз.

\cpic{1}{figures/2.6.png}

Тегістеуші түзуден $(1.3,5.4)$ және $(-0.7,-2.5)$ нүктелерін алайық. Онда көлбеулік коэффициенті
$$
n=\frac{(-2.5) - (+5.4)}{(-0.7) - (+1.3)}=3.95.
$$
Нәтиже \solref{14}-те алынған жауаптан сәл ғана үлкен.

\criterion[points=0.8]{Жауап $n\in[3.8,\ 4.1]$ аралығында}
\criterion[type=penalty, points=-0.4]{Жауап $n\in[3.5,\ 4.4]$ аралығында}
